\itemhead{p.58 练习 · 第2题（教材p.58）}{XJ-B1-C02-PR-P058-Y0432}
\label{XJ-B1-C02-PR-P058-Y0432}
{\footnotesize 湘教版教材p.58 / PDF第65页 / p.58 练习 / 第2题}
\field{原题}{行驶中的汽车，由于具有惯性，刹车后还要继续向前滑行一段距离才能停止，这段距离称为“刹车距离”。刹车距离是分析交通事故的一个重要指标。
\par 在一个限速为 $40\,\mathrm{km/h}$ 的弯道上，甲、乙两辆汽车相向而行，发现情况不对，同时刹车，但还是相碰了。事后现场勘查测得甲车的刹车距离略超过 $12\,\mathrm m$，乙车的刹车距离略超过 $10\,\mathrm m$，又知甲、乙两种车型的刹车距离 $s$（$\mathrm m$）与车速 $v$（$\mathrm{km/h}$）分别满足如下关系式：
\[s_{\text{甲}}=0.1v+0.01v^2,\qquad s_{\text{乙}}=0.05v+0.005v^2.\]
试用解一元二次不等式的方法判断甲、乙两辆汽车有无超速现象。}
\field{参考答案}{可严格推出 $v_{\text{甲}}>30$、$v_{\text{乙}}>40$，因此乙必超速；仅凭“略超过12米”无法严格判定甲是否超过40。若把“略超过”另解释为甲刹车距离仍小于20米，则甲未超速。答案册直接判甲未超速使用了这一未量化的语境假定，组卷时应补出上界。}
\field{核对说明}{原题条件不足：乙必超速；甲仅可推出v>30，要断言甲未超速需给出刹车距离上界。保持原题，不静默补条件。}
{\small \field{关联知识}{I04 一元二次不等式与参数分类（教材帮p.64、65、66 / PDF73、74、75）；F03 函数求值、解析式与函数建模（教材帮p.80、87、88、89、90、131 / PDF89、96、97、98、99、140）}
\field{方法与易错点}{化标准式、查系数、求根或判别式、按图定号，最后与实际范围取交集。 易错：套“内小外大”却忘开口；丢重根或降次；漏整数产量限制。}
\field{答案依据}{答案册相关题组：PDF第14页（册内页13）；按当前题库原题逐项重算，不按旧版题号机械对套。}}
